\subsubsection{Resultados esperados [Pendientes: 1. cerrar restricciones de datos y ejecución, 2. aprobar criterios de prueba]}\label{sec:inicio-resultados}
Inicio deberá producir una versión revisada de E1 con requisitos y trazabilidad, una arquitectura inicial E2 y un paquete E3 que documente los datos necesarios recolectados y preparados para el sector acordado, junto con las brechas que limiten su uso. La evidencia de cierre incluirá la reproducción del caso de práctica, la identificación de entradas y salidas, la revisión concreta de WFDS y las condiciones verificadas de acceso y ejecución, además de los criterios iniciales de E6. El equipo y el director deberán poder comprobar que el alcance corresponde a la información disponible y que existe una ruta técnica ejecutable hacia la construcción, por lo que una lista de fuentes sin preparación ni comprobación no bastará para acreditar ese resultado.

La contribución a OE1 será la definición del contrato de consulta y de las condiciones para seleccionar resultados, mientras que OE2 se apoyará en el catálogo previsto y en la viabilidad de ejecución. OE3 recibirá las configuraciones y entradas de práctica reproducibles, y OE4 contará con la arquitectura y las tareas de uso que orientarán su implementación. Estas relaciones no significan que los objetivos queden terminados en Inicio, pues sus productos se completarán mediante los incrementos posteriores. Cualquier restricción que impida sostener un resultado previsto deberá quedar documentada antes de considerar cerrada la fase.
